\documentclass[10pt,a4paper]{amsart}
\usepackage[margin=19mm]{geometry}
\usepackage{amsmath,amssymb,mathtools}
\usepackage[scr=rsfs,cal=euler]{mathalfa}
\usepackage{xfrac}
\usepackage[unicode,hidelinks,bookmarks=false]{hyperref}
\newcommand{\ie}{i.e.\ }
\newcommand{\cf}{cf.\ }
\newcommand{\resp}{respectively\ }
\newcommand{\Subsection}{\subsection}
\providecommand{\nobreakdash}{\nobreak\hbox{-}}
\setlength{\emergencystretch}{2em}
\multlinegap0pt
\usepackage{tikz}
\usepackage{enumerate}
\usepackage{xcolor}
\usepackage{xfrac}
\usepackage{bm}
\usepackage{mathtools}
\usepackage{csquotes}
\usepackage{amssymb}
\usepackage{stackengine}
\newtheorem{prop}{Proposition}
\newcommand{\HH}{\mathbb{H}}
\DeclareMathOperator{\Tr}{Tr}
\newcommand{\cL}{\mathcal{L}}
\newcommand{\sD}{\mathsf{D}}
\title{Performance Guarantees for Quantum Neural Estimation of Entropies}
\author{Sreejith Sreekumar, Ziv Goldfeld, Mark M. Wilde}
\date{Source: arXiv:2511.19289v2 (CC BY 4.0)}
\begin{document}
\maketitle

\noindent\emph{Source and licence.} Performance Guarantees for Quantum Neural Estimation of Entropies. Version arXiv:2511.19289v2 (CC BY 4.0), distributed by arXiv under the Creative Commons Attribution 4.0 International licence, http://creativecommons.org/licenses/by/4.0/, as recorded in the arXiv licence metadata for this exact version and shown on the abstract page https://arxiv.org/abs/2511.19289v2. Full licence notice: \texttt{licenses/arxiv-2511.19289-CC-BY-4.0.txt}.\par\smallskip

\noindent\textit{Editorial scope.} Counted unit: the article's prop eq:opt-omega-measure-rel-ent (source0.tex lines 1129--1141). The article's complete original proof of it is delivered with it (source0.tex lines 1142--1189, one \texttt{proof} environment of 48 lines). No mathematics is written, repaired or re-derived: the statement and the proof are the article's own lines, in the article's own notation, and no retained line is substituted or re-flowed.  Retained uncounted as the source-local context the proof needs: lines 318--321; lines 323--329; lines 1122--1127.  Nothing was omitted from any retained span, so the package carries no omission record.  The retained lines cite 3 works, whose entries are transcribed verbatim from the article's own bibliography source in the e-print and delivered with short editorial labels recorded in the item's reference mapping.  No retained line contains a figure environment, an image-inclusion command or drawing code, so no image is delivered and the package declares no asset dependency.  Editorial formatting only, in addition: an AMS \texttt{amsart} preamble that restates the article's own declarations and macros used by the retained lines, namely the environments \texttt{prop} on separate counters, together with the standard packages those lines need.  The retained environments print in this document as Proposition 1 in order of appearance, whereas the article prints the same items with the numbering of its own full document.  The complete native source of the article as distributed in the arXiv e-print for this exact version is bundled unchanged as \texttt{sources/189709/source0.tex}.
\begin{align}
\sD_{\mathsf{M}}(\rho\Vert\sigma)&=\sup_{\omega>0} \Tr[\rho\log\omega]+1-\Tr[\sigma\omega],\label{eq:measured-rel-ent-opt} \\
    \sD_{\mathsf{M},\alpha}(\rho\Vert\sigma)&= \frac{1}{\alpha-1} \log Q_{\alpha}(\rho\|\sigma), \quad \alpha \in (0,1) \cup (1,\infty), \label{eq:varformrenentmain2}
\end{align}

\begin{align}\label{eq:varformrenent2}
Q_{\alpha}(\rho\|\sigma)=\begin{cases}
\inf_{\omega>0}\big\{ \alpha\Tr[\rho\omega]+\left(1-\alpha\right)\Tr[\sigma\omega^{\frac{\alpha}{\alpha-1}}]\big\}, & \qquad\text{ for }\alpha  \in\left(0,\frac{1}{2}\right),\\
\inf_{\omega>0}\big\{ \alpha\Tr[\rho\omega^{1-\frac{1}{\alpha}}]+\left(1-\alpha\right)\Tr[\sigma\omega]\big\}, & \qquad\text{ for }\alpha  \in\left[\frac{1}{2},1\right),\\
\sup_{\omega>0}\big\{ \alpha\Tr[\rho\omega^{1-\frac{1}{\alpha}}]+\left(1-\alpha\right)\Tr[\sigma\omega]\big\}, & \qquad\text{ for }\alpha  \in\left(1,\infty\right).
\end{cases}
\end{align}

\begin{align} \label{eq:optrelvarform}
\omega_{\alpha}^{\star}(\rho,\sigma)= \begin{cases}
e^{(\alpha-1) H_{\alpha}^{\star}(\rho,\sigma)}=\sum_{i=1}^d\left(\frac{\Tr[P_{i,\alpha}^{\star} \rho]}{\Tr[P_{i,\alpha}^{\star} \sigma]} \right)^{\alpha-1} P_{i,\alpha}^{\star}, &\qquad\text{ for }\alpha  \in \left(0,\frac{1}{2}\right),  \\
e^{\alpha H_{\alpha}^{\star}(\rho,\sigma)}=\sum_{i=1}^d\left(\frac{\Tr[P_{i,\alpha}^{\star} \rho]}{\Tr[P_{i,\alpha}^{\star} \sigma]} \right)^{\alpha} P_{i,\alpha}^{\star}, &\qquad\text{ for }\alpha  \in\left[\frac{1}{2},\infty\right ), 
\end{cases}
\end{align}

The variational form in~\eqref{eq:varformrenentmain2} and the relation~\eqref{eq:optrelvarform} can be leveraged to show that $H_{\alpha}^{\star}(\rho,\sigma)$ is unique when $\rho,\sigma>0$, and should satisfy the  necessary (and sufficient) condition for optimality derived from it.  \begin{prop}[Necessary and sufficient condition for optimality]
Suppose $\rho,\sigma>0$. 
Then the optimal $\omega>0$ in~\eqref{eq:measured-rel-ent-opt} is unique and determined by the following equation:
\begin{align}
\sigma & =\int_{0}^{\infty}\ \left(\omega+sI\right)^{-1}\rho\left(\omega+sI\right)^{-1}ds.\label{eq:opt-omega-measure-rel-ent}
\end{align}
The
optimal $\omega>0$ in~\eqref{eq:varformrenent2} is unique and determined by the equations:
\begin{align}
\rho & =\left(\frac{1-\alpha}{\alpha}\right)\frac{\sin\!\left(\left(\frac{\alpha}{1-\alpha}\right)\pi\right)}{\pi}\int_{0}^{\infty}\ t^{\frac{\alpha}{\alpha-1}}\left(\omega+tI\right)^{-1}\sigma\left(\omega+tI\right)^{-1}dt,~~~ \forall~ \alpha\in\left(0,\frac{1}{2}\right),\label{eq:optimal-omega-alpha-0-half} \\
\sigma \mspace{-2 mu}& =\mspace{-2 mu}\left(\frac{\alpha}{1-\alpha}\right)\frac{\sin\!\left(\left(\frac{1-\alpha}{\alpha}\right)\pi\right)}{\pi}\mspace{-2 mu}\int_{0}^{\infty}\mspace{-2 mu} t^{\frac{\alpha-1}{\alpha}}\mspace{-2 mu}\left(\omega+tI\right)^{-1}\rho\left(\omega+tI\right)^{-1}dt, ~ \forall~ \alpha\in\left[\frac{1}{2},1\right)\cup\left(1,\infty\right). \label{eq:optimal-omega-alpha-half-infty}
\end{align}
\end{prop}

\begin{proof}
Let $A \in \cL(\HH_d)$, $\omega>0$ and $r\in\left(-1,0\right)\cup\left(0,1\right)$. We will use the following expressions  for the relevant matrix derivatives of interest  (see \cite{wilde2025-QIM} for a detailed derivation from first principles):  
\begin{align}
\frac{\partial}{\partial \omega}\Tr[A\log \omega] & =\int_{0}^{\infty}\ \left(\omega+sI\right)^{-1}A\left(\omega+sI\right)^{-1}ds,\label{eq:deriv-log}\\
\frac{\partial}{\partial \omega}\Tr[A\omega^{r}] & =\frac{\sin(r\pi)}{\pi}\int_{0}^{\infty}\ t^{r}\left(\omega+tI\right)^{-1}A\left(\omega+tI\right)^{-1}dt, \quad \forall r\in\left(-1,0\right)\cup\left(0,1\right).\label{eq:power-func-deriv-minus-1-to-plus-1}
\end{align}
The expressions in \eqref{eq:deriv-log} and \eqref{eq:power-func-deriv-minus-1-to-plus-1} also equal the  directional derivatives of $\log \omega$ and $\omega^r$, respectively, in the direction $A$ (see \cite[Proof of Theorem 1 and 5]{SB-IT-2025} for expressions of directional derivatives based on integral expressions given in \cite[Lemma 2.8]{Carlen-2010}). 


To determine  $\omega_{1}^{\star}(\rho,\sigma)$,
we set the matrix derivative  with respect to $\omega$ of the expression to be optimized in~\eqref{eq:measured-rel-ent-opt} to zero:
\begin{align}
0=\frac{\partial}{\partial\omega}\left(\Tr[\rho\log\omega]+1-\Tr[\sigma\omega]\right) & =\frac{\partial}{\partial\omega}\Tr[\rho\log\omega]-\frac{\partial}{\partial\omega}\Tr[\sigma\omega] \notag\\
 & =\int_{0}^{\infty}\ \left(\omega+sI\right)^{-1}\rho\left(\omega+sI\right)^{-1}ds-\sigma, \notag
\end{align}
where we applied~\eqref{eq:deriv-log}. We conclude that  $\omega_{1}^{\star}(\rho,\sigma)$ should satisfy~\eqref{eq:opt-omega-measure-rel-ent}. 
Uniqueness of $\omega_{1}^{\star}(\rho,\sigma)$ and sufficiency of the above condition for optimality follows from the strict concavity of the function $\omega\mapsto\Tr[\rho\log\omega]+1-\Tr[\sigma\omega]$
for $\rho,\sigma>0$,  which in turn follows from strict operator concavity of $\log $ and the fact that $\Tr[AB]> 0$ for $A,B >0$. 

To determine the optimal choice of $\omega$ in~\eqref{eq:varformrenent2} for $\alpha\in\left(0,\frac{1}{2}\right)$,
we set the  matrix derivative with respect to $\omega$ of the relevant expression to zero as follows:
\begin{align}
0&=  \frac{\partial}{\partial\omega}\left(\alpha\Tr[\rho\omega]+\left(1-\alpha\right)\Tr[\sigma\omega^{\frac{\alpha}{\alpha-1}}]\right) \nonumber \\
 &  =\alpha\rho+\left(\alpha-1\right)\frac{\sin\!\left(\left(\frac{\alpha}{1-\alpha}\right)\pi\right)}{\pi}\int_{0}^{\infty}\ t^{\frac{\alpha}{\alpha-1}}\left(\omega+tI\right)^{-1}\sigma\left(\omega+tI\right)^{-1}dt.\notag
\end{align}
Here, we used~\eqref{eq:power-func-deriv-minus-1-to-plus-1} with  $r=\frac{\alpha}{\alpha-1}\in\left(-1,0\right)$. Hence, $\omega_{\alpha}^{\star}(\rho,\sigma)$ satisfies~\eqref{eq:optimal-omega-alpha-0-half}.  Uniqueness of $\omega_{\alpha}^{\star}(\rho,\sigma)$ and sufficiency of the above equation for optimality follows from the strict convexity of the function $\omega\mapsto\alpha\Tr[\rho\omega]+\left(1-\alpha\right)\Tr[\sigma\omega^{\frac{\alpha}{\alpha-1}}]$
for $\rho,\sigma>0$, which in turn follows from strict operator convexity of $x \mapsto x^{r}$ for $r \in (-1,0)$ and $\Tr[AB]> 0$ for $A,B >0$.

To determine the optimal choice of $\omega$ in~\eqref{eq:varformrenent2} for $\alpha\in\left[\frac{1}{2},1\right)$, 
we set the following matrix derivative to zero:
\begin{align}
 0&= \frac{\partial}{\partial\omega}\left(\alpha\Tr[\rho\omega^{1-\frac{1}{\alpha}}]+\left(1-\alpha\right)\Tr[\sigma\omega]\right)\nonumber \\ 
  & =\left(1-\alpha\right)\sigma-\alpha\frac{\sin\!\left(\left(\frac{1}{\alpha}-1\right)\pi\right)}{\pi}\int_{0}^{\infty}\ t^{\frac{\alpha-1}{\alpha}}\left(\omega+tI\right)^{-1}\rho\left(\omega+tI\right)^{-1}dt. \notag 
\end{align}
Here, we applied~\eqref{eq:power-func-deriv-minus-1-to-plus-1} with $r=\frac{\alpha-1}{\alpha} \in (-1,0)$.
Hence, we conclude that $\omega_{\alpha}^{\star}(\rho,\sigma)$ should satisfy~\eqref{eq:optimal-omega-alpha-half-infty}. Uniqueness of $\omega_{\alpha}^{\star}(\rho,\sigma)$ and sufficiency of the condition above for optimality follows from the strict convexity of $\omega\mapsto\alpha\Tr[\rho\omega^{1-\frac{1}{\alpha}}]+\left(1-\alpha\right)\Tr[\sigma\omega]$, which is a consequence of strict operator convexity of $x \mapsto x^{r}$ for $r \in (-1,0)$ and $\Tr[AB]> 0$ for $A,B >0$.

Finally, note that $\omega_{\alpha}^{\star}(\rho,\sigma)$ for $\alpha \in (1,\infty)$ should satisfy 
\begin{align}
 0&= \frac{\partial}{\partial\omega}\left(\alpha\Tr[\rho\omega^{1-\frac{1}{\alpha}}]+\left(1-\alpha\right)\Tr[\sigma\omega]\right)\nonumber \\
  & =\alpha\frac{\sin\!\left(\left(\frac{\alpha-1}{\alpha}\right)\pi\right)}{\pi}\int_{0}^{\infty}\ t^{1-\frac{1}{\alpha}}\left(\omega+tI\right)^{-1}\rho\left(\omega+tI\right)^{-1}dt+\left(1-\alpha\right)\sigma, \notag
\end{align}
where we again applied~\eqref{eq:power-func-deriv-minus-1-to-plus-1} with $r=1-\frac{1}{\alpha}\in\left(0,1\right)$. Hence, 
 $\omega_{\alpha}^{\star}(\rho,\sigma)$ should satisfy~\eqref{eq:optimal-omega-alpha-half-infty}. 
Uniqueness of  $\omega_{\alpha}^{\star}(\rho,\sigma)$ and sufficiency of the above condition for characterizing $\omega_{\alpha}^{\star}(\rho,\sigma)$ follows from the strict concavity of the
function $\omega\mapsto\alpha\Tr[\rho\omega^{1-\frac{1}{\alpha}}]+\left(1-\alpha\right)\Tr[\sigma\omega]$
for $\rho,\sigma>0$, which in turn follows from strict operator concavity of  $x \mapsto  x^r$ for $r \in (0,1)$. 
\end{proof}

\begin{thebibliography}{99}
\bibitem[Carlen]{Carlen-2010} Eric~A. Carlen. \newblock ``Trace inequalities and quantum entropy: An introductory course''. \newblock In Contemporary Mathematics. \newblock Volume 529 of Entropy and the Quantum. \newblock American Mathematical Society~(2010).
\bibitem[SB-IT-]{SB-IT-2025} Sreejith Sreekumar and Mario Berta. \newblock ``Limit distribution theory for quantum divergences''. \newblock \href{https://dx.doi.org/10.1109/TIT.2024.3482168}{IEEE Transactions on Information Theory {\bf 71}, 459--484}~(2025).
\bibitem[wilde2]{wilde2025-QIM} Mark~M. Wilde. \newblock ``{Quantum Fisher information matrices from R\'enyi relative entropies}''~(2025). \newblock \href{http://arxiv.org/abs/2510.02218}{arXiv:2510.02218}.
\end{thebibliography}
\end{document}
